\section{The instruction set}\label{sec:isa}
% =====================================================================================

\subsection{The idea of a Pauli frame}

\begin{plain}
Here is the trick that makes the machine cheap.

A quantum state is expensive to write down: for two qutrits you need nine complex
numbers, and keeping them accurate is most of what makes quantum computing hard. But for
a large and useful family of operations --- the \emph{Clifford} operations --- you never
need the state at all. You only need a \emph{label} saying how the state has been shifted
relative to where it started.

That label is four trits. Four trits is eight bits. The machine's entire ``register
file'' for tracking Clifford operations is one byte.
\end{plain}

\headline{The frame unit tracks a label, not a state. Four trits, $81$ possible values,
eight bits of memory --- and it correctly accounts for a group of $4{,}199{,}040$
distinct operations.}

\begin{spec}
The Pauli frame is $(x_p,z_p,x_f,z_f)\in\mathbb{F}_3^4$, denoting the operator
$X^{x_p}Z^{z_p}\otimes X^{x_f}Z^{z_f}$. A Clifford unitary $U$ acts on frames by
conjugation, $U(X^aZ^b)U^\dagger = X^{a'}Z^{b'}$ up to phase, and the induced map on
$\mathbb{F}_3^4$ is \emph{linear and symplectic}: it preserves
\[
  \langle u,v\rangle \;=\; (x_p z_p' - z_p x_p') + (x_f z_f' - z_f x_f') \bmod 3 .
\]
\end{spec}

\subsection{The eight opcodes}

\begin{spec}[$\mathcal{I}_{\mathrm{holo}}$, complete semantics]
\begin{center}\small
\begin{tabular}{@{}llll@{}}
\toprule
\textbf{Op} & \textbf{Name} & \textbf{Action on $(x_p,z_p,x_f,z_f)$} & \textbf{Kind}\\
\midrule
\texttt{000} & $F_p$ & $(x_p,z_p)\mapsto(-z_p,\;x_p)$              & linear \\
\texttt{001} & $F_f$ & $(x_f,z_f)\mapsto(-z_f,\;x_f)$              & linear \\
\texttt{010} & $S_p$ & $z_p\mapsto z_p+x_p$                        & linear \\
\texttt{011} & $S_f$ & $z_f\mapsto z_f+x_f$                        & linear \\
\texttt{100} & $\mathrm{CX}$ & $d{=}0$: $z_p\mapsto z_p-z_f$, $x_f\mapsto x_f+x_p$
                              & linear \\
             &               & $d{=}1$: $x_p\mapsto x_p+x_f$, $z_f\mapsto z_f-z_p$
                              & linear \\
\texttt{101} & $\sigma^5{=}Z$ & $z_p\mapsto z_p+1$ \ or\  $z_f\mapsto z_f+1$
                              & \textbf{affine} \\
\texttt{110} & $D_{12}$-mirror & transport in the dihedral group $D_{12}$ & (not a frame op)\\
\texttt{111} & $M_{36}$-magic & typed request for one of $36$ Witting rays & (handshake)\\
\bottomrule
\end{tabular}
\end{center}
\end{spec}

\subsection{The reachability theorem, and why it matters}

\begin{plain}
Now a question that sounds trivial and is not: \emph{starting from a blank machine, which
states can these instructions actually reach?}

The answer is startling. All six of the Clifford instructions, applied in any order, any
number of times, reach exactly \emph{one} state --- the blank one. They cannot move the
machine at all.

The reason is simple once you see it. Those six instructions are all \emph{linear}, and a
linear map always leaves the origin where it is. Multiplying zero by anything gives zero.
So a machine built from only those six instructions is, quite literally, a machine that
does nothing, no matter how sophisticated its circuitry looks.

The seventh instruction, $\sigma^5=Z$, is different: it \emph{adds one}. Addition moves
the origin. That single instruction is what makes all $81$ states reachable --- and hence
what makes the machine a computer instead of an ornament.
\end{plain}

\begin{center}
\begin{tikzpicture}[font=\small,>={Latex[length=2.4mm]}]
  % left: clifford only
  \begin{scope}
    \draw[draw=rule,rounded corners] (-0.6,-1.5) rectangle (3.6,1.9);
    \node[font=\bfseries\small] at (1.5,1.55) {Clifford opcodes only};
    \foreach \x in {0,...,4}\foreach \y in {0,...,3}{
      \node[circle,draw=rule,inner sep=1.1pt,fill=white] at (0.7*\x,0.55*\y-1.0) {};}
    \node[circle,fill=bad,inner sep=2.6pt] at (0,-1.0) {};
    \draw[->,line width=0.9pt,draw=bad] (0.05,-0.85) arc (-40:250:0.28);
    \node[font=\scriptsize,text=bad,align=center] at (1.9,-1.35)
      {reachable: \textbf{1} of $81$};
  \end{scope}
  % right: with Z
  \begin{scope}[xshift=7.2cm]
    \draw[draw=rule,rounded corners] (-0.6,-1.5) rectangle (3.6,1.9);
    \node[font=\bfseries\small] at (1.5,1.55) {with $\sigma^5=Z$};
    \foreach \x in {0,...,4}\foreach \y in {0,...,3}{
      \node[circle,fill=good,inner sep=2.0pt] at (0.7*\x,0.55*\y-1.0) {};}
    \node[font=\scriptsize,text=good,align=center] at (1.9,-1.35)
      {reachable: \textbf{81} of $81$};
  \end{scope}
  \draw[->,line width=1.1pt,draw=ink] (4.1,0.2) -- (6.3,0.2)
    node[midway,above,font=\scriptsize,align=center] {one\\instruction};
\end{tikzpicture}
\end{center}

\begin{spec}[Proved by exhaustive search, Pass 2774]
Breadth-first search over all $3^4=81$ frames from $(0,0,0,0)$:
\[
\begin{array}{lr}
\textrm{$\mathrm{CX}$ alone} & 1 \\
\textrm{all six symplectic opcodes} & 1 \\
\textrm{full ISA (adding $\sigma^5{=}Z$)} & 81 \\
\end{array}
\]
Reproduce: \cmd{py -3 analysis/w33\_pass2774\_isa\_frame\_reachability.py}\\
Certificate: \cmd{data/PART\_W33\_PASS2774\_ISA\_FRAME\_REACHABILITY.json}
\end{spec}

\subsection{How much computation the eight opcodes actually buy}

\begin{plain}
``All $81$ states are reachable'' is weaker than it sounds --- it says you can get
anywhere, not that you can perform every operation. The stronger question is: what is the
complete list of things this instruction set can do?
\end{plain}

\begin{spec}[Proved, Pass 2778]
The six Clifford opcodes generate a group of order exactly $51{,}840$ --- that is, the
whole of $\mathrm{Sp}(4,3)$, nothing missing. Adjoining the translation gives
\[
  \langle \mathcal{I}_{\mathrm{holo}}\rangle
   \;=\; \mathrm{ASp}(4,3)\;=\;\mathbb{F}_3^4\rtimes\mathrm{Sp}(4,3),
  \qquad |\cdot| = 81\times 51840 = 4{,}199{,}040 .
\]
\textbf{One translation generator suffices.} Enabling only $Z$ on the past register gives
the same group as enabling both, because $\mathrm{Sp}(4,3)$ is transitive on the $80$
nonzero frame vectors: the orbit of a single translation already spans $\mathbb{F}_3^4$.
Measured orbit size: $80$. A second translation instruction adds nothing and can be
removed from the ISA.
\end{spec}

\headline{Eight opcodes. Three bits. $4{,}199{,}040$ distinct operations, exactly ---
and the eighth opcode is provably redundant.}

\subsection{How small can the instruction set be?}

\begin{plain}
If one opcode is redundant, how many are? This is not idle: every instruction removed is
decoder logic removed, and if enough go, the opcode field itself gets shorter --- which
shortens every instruction word in the machine.

The answer is that you can throw away half of them.
\end{plain}

\begin{spec}[Exhaustive over all subsets, Pass 2789]
Because $\mathbb{F}_3^4\rtimes H=\mathrm{ASp}(4,3)$ iff $H=\mathrm{Sp}(4,3)$ and at least
one translation is present, the search reduces to the six linear opcodes:
\[
\begin{array}{ll}
\text{subsets of size }1\text{ generating }\mathrm{Sp}(4,3): & 0\\
\text{subsets of size }2: & 0\\
\text{subsets of size }3: & \mathbf{6}
\end{array}
\]
The six minimal triples are
$\{F_p,F_f,\mathrm{CX}_{p\to f}\}$, $\{F_p,F_f,\mathrm{CX}_{f\to p}\}$,
$\{F_p,S_f,\mathrm{CX}_{p\to f}\}$, $\{F_p,\mathrm{CX}_{p\to f},\mathrm{CX}_{f\to p}\}$,
$\{F_f,S_p,\mathrm{CX}_{f\to p}\}$, $\{F_f,\mathrm{CX}_{p\to f},\mathrm{CX}_{f\to p}\}$.
\end{spec}

\headline{Three Clifford opcodes plus one translation generate everything.
Four frame instructions --- a \textbf{two-bit} opcode field, not three.}

\begin{plain}
Two details are worth noticing. First, there are \emph{six} different minimal choices, so
a chip designer has real freedom to pick whichever three are cheapest in silicon rather
than being handed one answer.

Second, every single one of the six contains at least one Fourier gate \emph{and} at
least one controlled-add. No collection of phase gates and one Fourier gate will do. The
entangler is not optional --- which is the formal version of the intuition that a machine
that never couples past to future is not really using the self-entanglement at all.
\end{plain}

% =====================================================================================
\section{The hardware}\label{sec:hardware}
% =====================================================================================

\subsection{The frame unit datapath}

\begin{center}
\begin{tikzpicture}[font=\footnotesize,>={Latex[length=2.2mm]},
  reg/.style={draw=ink,fill=specbg,minimum width=13mm,minimum height=8mm,rounded corners=1.5pt},
  alu/.style={draw=ink,fill=plainbg,minimum width=17mm,minimum height=8mm,rounded corners=1.5pt},
  mux/.style={draw=ink,fill=white,trapezium,trapezium angle=68,minimum width=9mm,
              minimum height=7mm,shape border rotate=270,inner sep=1pt},
  w/.style={-{Latex[length=1.8mm]},draw=rule,line width=0.7pt}]

  % registers
  \node[reg] (xp) at (0,3.0)  {$x_p$};
  \node[reg] (zp) at (0,1.9)  {$z_p$};
  \node[reg] (xf) at (0,0.8)  {$x_f$};
  \node[reg] (zf) at (0,-0.3) {$z_f$};
  \node[font=\scriptsize\itshape,text=rule,above=0.5mm of xp] {$4$ trits $=$ $8$ bits of state};

  % f3 units
  \node[alu] (a1) at (3.6,2.6) {$\mathrm{add}_3$};
  \node[alu] (a2) at (3.6,1.4) {$\mathrm{sub}_3$};
  \node[alu] (a3) at (3.6,0.2) {$\mathrm{neg}_3$};
  \node[font=\scriptsize\itshape,text=rule,align=center] at (3.6,-0.9)
    {mod-$3$ arithmetic\\(pure LUT logic, no carry chain)};

  % opcode decode
  \node[alu,minimum height=22mm,minimum width=15mm,fill=future!12] (dec) at (7.2,1.4)
    {\begin{tabular}{c}opcode\\decode\\\scriptsize $3$ bits\end{tabular}};

  % next-state mux
  \node[mux] (m) at (10.0,1.4) {};
  \node[font=\scriptsize] at (10.0,1.4) {mux};

  \node[reg,fill=warnbg] (ld) at (10.0,3.4) {load port};
  \node[font=\scriptsize\itshape,text=rule,align=center,above=0.5mm of ld]
    {\textbf{without this the whole unit}\\\textbf{synthesises to nothing}};

  \foreach \r in {xp,zp,xf,zf}{
    \draw[w] (\r.east) -- ++(0.7,0) |- (a1.west);
    \draw[w] (\r.east) -- ++(0.7,0) |- (a2.west);
  }
  \draw[w] (a1.east) -- (dec.west|-a1);
  \draw[w] (a2.east) -- (dec.west|-a2);
  \draw[w] (a3.east) -- (dec.west|-a3);
  \draw[w] (dec.east) -- (m.west);
  \draw[w] (ld.south) -- (m.north);
  \draw[w] (m.south) |- ++(0,-1.4) -| (-1.4,-1.4) |- (xp.west);
  \draw[w] (-1.4,-1.4) |- (zp.west);
  \draw[w] (-1.4,-1.4) |- (xf.west);
  \draw[w] (-1.4,-1.4) |- (zf.west);

  \node[font=\scriptsize,text=rule] at (5.6,-1.6)
    {\cmd{rtl/w33\_pass2777\_isa\_frame\_unit.sv}};
\end{tikzpicture}
